\documentclass[	pdftex, 
								a4paper,
								11pt, DIV11, BCOR5mm,
								parskip,
								%openany,
								]{scrreprt}
\usepackage[utf8]{inputenc}
\usepackage[T1]{fontenc}

\usepackage[colorlinks=true,linkcolor=blue,citecolor=blue]{hyperref}
\usepackage{comment}
\usepackage[ngerman]{babel}
\usepackage[nonumberlist,toc,acronym]{glossaries}

\usepackage{fancyhdr}
\usepackage[section]{placeins}
\usepackage{graphicx}
\usepackage{listings}
\usepackage[german]{babelbib}

\usepackage{varwidth}
\usepackage{float}

\usepackage{epstopdf}
\usepackage{textpos}
\usepackage{changepage}
\usepackage{theorem}
\usepackage{caption}
\usepackage{framed}
\usepackage{booktabs}
\usepackage{tikz}
\usetikzlibrary{decorations.markings,arrows,positioning,trees,calc,fit,shapes}
\usepackage[noend]{algorithmic}
\usepackage{algorithm}
\usepackage{environ}
\usepackage{ulem}

%Kopfzeile hinzufuegen
\renewcommand{\sectionmark}[1]{\markboth{#1}{}} % set the \leftmark

\fancypagestyle{plain}{
    
    \fancyhead[R]{\leftmark}
    \fancyhead[L]{\myauthor}
    \renewcommand{\headrulewidth}{1pt}
    %\fancyfoot[R]{\thepage}
}

\pagestyle{fancy}

%Literaturverzeichnis hinzufuegen
\bibliographystyle{babplain}

% Glossar hinzufuegen
\makeglossaries
\newacronym{usda} {USDA}{US Department of Agriculture}
\newacronym{fdc}{FDC}{FoodData Central}
\newacronym{rest}{REST}{Representational state transfer}
\newacronym{dge}{DGE}{Deutsche Gesellschaft für Ernährung}
\newacronym{pal}{PAL-Value}{physical activity level-Value}
\newacronym{json}{JSON}{JavaScript Object Notation}
\newacronym{iri}{IRI}{International Resource Identifier}
\newacronym{http}{HTTP}{HyperText Transfer Protocol}
\newacronym{xml}{XML}{Extensible Markup Language}
\newacronym{cdno}{CDNO}{Compositional Dietary Nutrition Ontology}
\newacronym{obo}{OBO}{Open Biological and Biomedical Ontology}
\newacronym{sparql}{SPARQL}{SPARQL Protocol And RDF Query Language}
\newacronym{sql}{SQL}{Structured Query Language}
\newacronym{uri}{URI}{Uniform Resource Identifier}
\newacronym{gi}{GI}{Glykämischer Index}
\newacronym{bmi}{BMI}{Body-Mass-Index}



\begin{document}
\setcounter{page}{-2}
\pagestyle{empty}
\pagenumbering{none}
\newcommand{\firstreviewer}{Prof. Michael Beetz PhD}
\newcommand{\secondreviewer}{}
\newcommand{\supervisor}{Michaela Kümpel}
\newcommand{\thesistype}{Master Thesis}
\newcommand{\myauthor}{Naser Azizi, Sorin Arion}
\newcommand{\mymaintitle}{A Knowledgebase with which you can generate robot plan for multiple mixing actions}
\newcommand{\mysubtitle}{A nutrition ontology with web interface for customer specific dietary \newline recommendations} 
\newcommand{\mytitle}{\centering {\Huge \mymaintitle}\\[.3in] 
    {\Large \mysubtitle}}
\newcommand{\pdftitle}{\mymaintitle - \mysubtitle}
\newcommand{\formattedfronttitle}{{\mytitle}}
\newcommand{\formattedinnertitle}{{\mytitle}}

\begin{titlepage}
	\vspace*{-2.2cm}
	\begin{adjustwidth}{-0cm}{-2.3cm}
	\thispagestyle{empty}
        \begin{figure}
            \begin{minipage}{.4\linewidth}
	\begin{flushleft}
		\includegraphics[height=1.5cm]{Graphics/unilogo-transp.pdf}
	\end{flushleft}
    \end{minipage}
    \hspace{.2\linewidth}
            \begin{minipage}{.4\linewidth}
	\begin{flushright}
		\includegraphics[height=3.0cm]{Graphics/logo-ai-small.pdf}
	\end{flushright}
    \end{minipage}
\end{figure}
	  \vfill
	  %\scalebox{0.95}{
    	  %\begin{minipage}{1.2\textwidth}
    	  %{
    		%\formattedfronttitle 
    	  %}
    	  %\end{minipage}
	  %}
	\begin{center}
	  {\huge \mymaintitle}\\
	  
	  \vfill
	  {\Large \thesistype}\\[2.5ex]
	  {\Large\em \myauthor}
	  \vfill
	{  
      \renewcommand\arraystretch{1.5}
      \begin{tabular}{l@{\hspace{2em}}r@{\hspace{1ex}}p{7cm}}
     Pr\"ufer der \thesistype: & 1. & \firstreviewer\\
                                 & 2. & \secondreviewer\\
	Supervisor		    &   & \supervisor\\
   \end{tabular}
  }
	\end{center}
	\end{adjustwidth}
\end{titlepage}
	\pagenumbering{arabic}
	\setcounter{page}{1}
	\pagestyle{plain}
	%\addcontentsline{toc}{chapter}{Eidesstaatliche Erkl\"arung}
	\chapter*{Eidesstattliche Erkl\"arung}
	
	Hiermit erkl\"aren wir, dass die vorliegende Arbeit selbstst\"andig angefertigt,
	nicht anderweitig zu Pr\"ufungszwecken vorgelegt und keine anderen als die
	angegebenen Hilfsmittel verwendet habe. S\"amtliche wissentlich verwendete
	Textausschnitte, Zitate oder Inhalte anderer Verfasser wurden ausdr\"ucklich als
	solche gekennzeichnet.
	
	Bremen, den \makeatletter\@date\makeatother
	
	\vspace*{1em}
	\rule{15em}{0.16667pt}\\
	\author{Naser Azizi, Sorin Arion}
	\makeatletter\@author\makeatother
	
	
	\normalem
	%Abstrakt
	%\addcontentsline{toc}{chapter}{Abstract}


	\chapter*{Motivation}
	
	Although industrial robots have been considered standard for many years, the widespread adoption of AI-based autonomous household robots is still far from being anticipated in every household.
	Obstacles such as unfamiliar surroundings and a limited knowledge base can significantly hinder the ability of an autonomous robot to effectively respond in a given situation, hence, the robot must precisely plan for every minor detail, even those that might appear trivial to us humans.
	To achieve the defined objectives, a proficient and comprehensive robotic system including Control, Perception, Navigation, and Knowledge is necessary.
	An illustrative goal could be cooking. To execute this task, the robot must possess various capabilities, including placing items, grasping objects, pouring liquids, cutting ingredients, and mixing components.	Those things seem trivial to us humans, but for the robots it requires a lot of Implementation and information.
	In this thesis, we aim to introduce an approach detailing how a robotic system can execute mixing actions effectively.	
	The challenges arise from the diverse methods of mixing, which further vary based on the specific ingredients being combined. Additionally, the consideration of compatible Mixing Tools and Mixing Containers is crucial, as not every tool can be utilized with every container.
	Through the incorporation of a knowledge graph containing rules related to to various actions, we aim to empower the robotic system to make informed decisions on the appropriate motions to employ. This decision-making process will take into account the specific task at hand and the involved ingredients.
	By incorporating this knowledge graph, we aim to advance towards autonomous robots capable of engaging in cooking activities.

	\chapter*{Introduction}
	First, we intend to analyze existing research, particularly examining studies that delve into diverse mixing motions and exploring related systems that contain a queryable knowledge graph.
	In the following section, we will introduce the reader with the frameworks employed to accomplish our objectives.	
	Subsequently, we will present our approach to data acquisition and data representation.
	In this chapter, we aim to explain the mixing tasks under consideration and articulate our methodology for representing this data to enable effective querying.
	In the subsequent section, we will illustrate the implementation of our rules, with the ultimate aim of determining the appropriate motion to be employed.
	To validate our concept, we have opted to simulate various scenarios and assess their outcomes, demonstrating the efficacy of our implemented system.
	Finally, we will delve into the discussion of our results and draw conclusions to wrap up our work.
	
	\chapter*{Related Work}
	\begin{itemize}
		\item https://fluidlab2023.github.io/
		\item 
	\end{itemize}

	\chapter*{How does the robot works / plans, parameters, used tools, etc}
	In this chapter we want to introduce the reader to the frameworks and libraries we are using, also we want to briefly present the robot that should be able to perform the implemented tasks.
	\section*{Pr2}
	Introduced in 2010 by Willow Garage (https://robotsguide.com/robots/pr2), the PR2 stands as an advanced research robot. Boasting multiple joints and 20 degrees of freedom, this robot excels in autonomous navigation and the manipulation of a diverse array of objects, making it an ideal choice for our specific needs. Additionally, it is equipped with a HeadStereoCamera that can be used to perceive the surroundings.
	\section*{PyCram}
	Robotic systems are complex, comprising various components such as Knowledgebase, Perception, Manipulation/Control, Navigation, and more. Managing these diverse elements can pose challenges and requires multiple interfaces between them. One potential strategy is to integrate all these interfaces into a unified framework. An example of such a comprehensive framework is PyCram(https://github.com/cram2/pycram).
	Implemented in Python3, the PyCram framework incorporates ROS for robot communication. Utilizing PyCram, we can execute High-Level plans on the PR2 to evaluate our implementation.
	\section*{Ontology}
	Ontologies are structured frameworks that provide a formal representation of knowledge within a specific domain. They play a crucial role in knowledge representation, facilitating the organization and sharing of information in a way that is both machine-readable and understandable by humans. 

	Key components of ontologies include:

	\begin{itemize}
		\item Concepts/Classes: These represent abstract or concrete entities within a domain. For example, in a medical ontology, "Patient" and "Disease" might be classes.
		\item Properties/Roles: These define the relationships between concepts. For instance, in a social network ontology, "FriendOf" could be a property connecting individuals.
		\item Instances/Individuals: These are specific members or examples of a class. In a geographical ontology, "New York City" and "Paris" could be instances of the class "City."
		\item Axioms: These are statements that describe the properties and relationships of the entities within the ontology. Axioms help define the logic and rules governing the domain.
		\item Hierarchy: Ontologies often organize concepts into a hierarchy, with more general concepts at the top and more specific ones below. This hierarchical structure aids in categorization and understanding.
		\item Inference Rules: These rules define how new information can be derived from existing information in the ontology. Inferences help systems reason and make deductions based on the knowledge encoded in the ontology.
	\end{itemize}
	We utilize ontologies in knowledge representation and reasoning systems to empower the robot with the ability to comprehend and handle information in a structured fashion for our specific objectives.
	\section*{OWLReady}
	One important library used for our Implementation is OWLReady. "OWLReady" is a Python library designed for ontology-oriented programming. It facilitates the development, manipulation, and querying of ontologies using the Web Ontology Language (OWL), a standard for representing knowledge in a machine-readable format. OWLReady simplifies ontology-related tasks by providing a convenient and object-oriented interface for working with OWL ontologies in Python.

	Key features of the OWLReady library include:

	\begin{itemize}
		\item Object-Oriented Programming (OOP): OWLReady adopts an object-oriented approach, allowing users to interact with ontology entities as Python objects. This makes it more intuitive for developers familiar with Python's OOP principles.
		\item Ontology Loading and Parsing: The library supports the loading and parsing of OWL ontologies, making it easy to access and manipulate ontology data within Python scripts or applications.
		\item Class and Individual Manipulation: OWLReady provides functionality for creating, modifying, and querying classes and individuals within an ontology. This allows for dynamic and programmatic management of ontology content.
		\item Reasoning Support: Depending on the version and features, OWLReady may offer support for reasoning tasks. Reasoning involves deducing implicit information based on the logical relationships defined in the ontology.
		\item Integration with RDFLib: OWLReady may integrate with RDFLib, another Python library commonly used for working with Resource Description Framework (RDF) data. This integration enhances the capabilities of handling semantic data.
	\end{itemize}

	\section*{SWRL}
	SWRL, which stands for Semantic Web Rule Language, is a rule language that allows users to define rules about the relationships between classes and individuals in ontologies represented in the Web Ontology Language (OWL). SWRL is part of the W3C's Semantic Web technology stack and is designed to be used in conjunction with OWL to express complex relationships and infer new information based on existing knowledge.

	SWRL Rules have a specific syntax and consist of two main components:
	\begin{itemize}
		\item Antecedent (Body): This part of the rule specifies the conditions or constraints that must be satisfied for the rule to be applicable. It describes the current state of the ontology that triggers the rule.
		\item Consequent (Head): This part defines the actions or inferences that should be taken if the conditions specified in the antecedent are satisfied. It describes the changes or additional information that should be inferred when the rule is triggered.
	\end{itemize}
	SWRL supports various built-in predicates and functions, and users can create their own custom rules to suit their specific ontology. Some common elements in SWRL rules include:
	\begin{itemize}
		\item Individuals: Refers to specific instances of classes in the ontology.
		\item Class and Property Relationships: Describes relationships between classes and properties in the ontology.
		\item Built-in Predicates and Functions: Includes operations such as arithmetic, string manipulation, and comparison functions that can be used in the rule conditions.
	\end{itemize}
	Here's a simple example of a SWRL rule:
	\newline
	$Person(?x) ^ hasChild(?x, ?y) -> Grandparent(?x, ?z) ^ hasChild(?y, ?z)$
	\newline
	In this example:

	If an individual (?x) is a Person and has a child (?y),
	
	Then, infer that the individual (?x) is a Grandparent of another individual (?z), and the child (?y) is the parent of (?z).
	
	SWRL rules are useful for expressing complex relationships and constraints within ontologies, enabling automated reasoning systems to make inferences and derive new knowledge from existing data.

	\chapter*{Data acquisition}
	\begin{itemize}
		\item Video analysis
		\item Structured the required parameters for the knowledge base
	\end{itemize}
	The first step in acquiring the needed data, was to acknowledge which task varations of mixing are actually important. So we had to analyze the word "Mixing" and its hyponyms. 
	hyponyms are subordered words of a given word, for example one hyponym of mixing could be beating. By looking on different websites like "framenet, Wordnet" , we created a list of possible hyponyms, that could be taken under consideration.
	For all those hyponyms, we delegated a WikiHow extraction search which should show us, how many times one of these words occur, in the context of cooking.
	In the table below the results can be seen.
	
	HIER TABELLE

	From this table one can see that many hyponyms have a low occurance and will be therefore not considered.
	We consider the words: Mixing, Stiring, Beating, Combining, Whiping, Whisking and Folding. Though Blending achieved a moderate occurance, we decided upon not considering it, because it involves a lot of electronic mixing machines, which the robot could not handle.
	From now on we will refer to these words also as Tasks.
	After deciding upon which tasks we consider, we have to analyze every task execution and define a structure.
	We started by analyzing different videos on WikiHow that included at least one of our choosen task and divided the information by:
	
	Task, used Tools, used Container, Ingredients, Motion.

	\begin{itemize}
		\item Tasks: The tasks can differ in its execution, like the tasks Beating (Mixing can be gentle or vigorous, while beating involves a more forceful and rapid action) and Folding (To maintain a light and fluffy texture in a cake, folding is often used to add dry ingredients without overmixing), which will ultimately lead to a decision upon which Motion should be used in the specific use case.
		\item Tools: Different Tools are used for different tasks, which also has a correlation with the choosen Container.
		\item Container: Most Mixing Tasks use a mixing bowl for container, but also container like Pot and Pan when considering cooking in boiling water.
		\item Ingredients: Besides Tasks, the ingredients are the most important component of our knowledgebase regarding the decision making upon motions. We divide the Ingredients in 4 subcategories: Dry/Powder, Wet, Liquid and Solid.
		\item Motions: The motions used in the analyzed videos will help us upon creating a knowledgebase from which the robot should be able to infer which motions he will use in different situations.
	\end{itemize}

	In the following we want to present our analysis results regarding the videos from WikiHow, but also from well known cooks, like Jamie Oliver:

	HIER TABELLE

	From this videos we extracted informations about the executed motions. The following motions were defined:
	\begin{itemize}
		\item Circular
		\item Whirlstorm
		\item Folding
		\item Diagonal
		\item etc
	\end{itemize}

	By analyzing the videos we concluded that the motion decision is based on the task and Ingredients. Upon this thoughts we decided to design a decision tree regarding how the motions will be choosen by the robot.

	HIER DECISION tree


	\chapter*{Data representation}
	\begin{itemize}
		\item Ingredients, Tools, etc.
		\item Different Tasks/Actions maps on different Motions
		\item Motion dependencies
		\item Fobidden Tools with Containers
		\item How is an action defined
		\item Examples of using the knowledge graph
	\end{itemize}
	\chapter*{Implementation}
	\begin{itemize}
		\item OWLReady
		\item SWRL
		\item Inferring parameters for the actual plan
	\end{itemize}
	\chapter*{Simulation}
	\begin{itemize}
		\item BulletWorld as Simulation platform -> how is the robot actually shown, which limitations exist
		\item Quering the knowledge base
		\item Inferring the parameters in the Simulation
		\item Task execution with multiple sequential Motions
	\end{itemize}
	\section*{Simulation to RealWorld gap}
	\begin{itemize}
		\item Big Problem: Uncertaininty
		\item Perception Solution as first approach: RoboKudo
		\item Data acquisition: Blenderproc
		\item Model Training: YoloV8
		\item Results showing the real world Perception.
	\end{itemize}
	\chapter*{Evaluation}
	\begin{itemize}
		\item Evaluation of multiple motions multiple times
		\item Results of that
	\end{itemize}
	\chapter*{Summary / Fazit}
avc
\end{document}

